\section{Limitaciones actuales}
\label{ch:limitaciones}
\label{sec:etapas_pendientes}

El presente informe documenta el estado del proyecto al 13 de octubre de 2026. Las siguientes limitaciones son \textbf{parte inherente del estado actual de la investigación}, no fallas del proyecto:

\subsection{Auditoría estadística integral pendiente}
\label{sec:auditoria_pendiente_limitaciones}

Los diagnósticos reportados en el Capítulo \ref{ch:diagnostico} (KMO $\approx 0,642$, Bartlett significativo, MSA estrato 1--2 $\approx 0,414$, 32 municipios con $|z|>4$, sensibilidad NR) son \textbf{preliminares}. Falta:

\begin{itemize}
    \item Comparación sistemática base total vs. base válida (correlaciones, KMO, MSA).
    \item Evaluación de transformaciones para asimetría/curtosis (ej. techo en gas, internet).
    \item \textit{Parallel analysis} para retención de componentes.
    \item Análisis de estabilidad (bootstrap, sensibilidad a outliers).
    \item Decisión fundamentada sobre variables problemáticas (estrato 1--2, hogares con menores, dirección adultos 60+).
\end{itemize}

\textbf{No se ha aprobado la matriz para PCA definitivo.}

\subsection{PCA no ejecutado}
\label{sec:pca_pendiente}

\begin{itemize}
    \item No se han estimado componentes principales.
    \item No se han inspeccionado eigenvalues, scree plot, cargas, varianza explicada.
    \item No se ha evaluado interpretabilidad del PC1 como ``vulnerabilidad general''.
    \item No se ha definido criterio de retención (Kaiser, scree, parallel analysis, varianza acumulada).
\end{itemize}

\subsection{IVS (Índice de Vulnerabilidad Socioeconómica) no construido}
\label{sec:ivs_pendiente}

\begin{itemize}
    \item No existe un índice sintético validado.
    \item No se han calculado scores municipales de vulnerabilidad.
    \item No se ha orientado el signo (score alto = mayor vulnerabilidad).
    \item No se ha evaluado estabilidad del ranking municipal.
\end{itemize}

\subsection{Validación externa con IPM pendiente}
\label{sec:validacion_ipm_pendiente}

\begin{itemize}
    \item El IPM Municipal agregado está reservado como criterio de validación.
    \item No se ha correlacionado IVS con IPM (evitando circularidad por exclusión de indicadores redundantes).
    \item No se ha analizado concordancia de rankings, análisis de residuos, ni sensibilidad de la validación.
\end{itemize}

\subsection{Integración SECOP II pendiente}
\label{sec:secop_pendiente}

\begin{itemize}
    \item No se han obtenido ni integrado datos de SECOP II.
    \item No se ha definido ventana temporal definitiva (exploratoria 2015--2025).
    \item No se han construido indicadores municipales-año de contratación (valor, número, concentración, sectores).
    \item No se ha resuelto el \textit{crosswalk} territorial definitivo (códigos 94663, 27493 pendientes de verificación DANE/IGAC).
    \item No se ha definido denominador para tasas de contratación (PIB municipal, población, presupuesto, referencia nacional).
\end{itemize}

\subsection{Análisis de asociación/correspondencia pendiente}
\label{sec:asociacion_pendiente}

\begin{itemize}
    \item Condicionado a la existencia de IVS validado y datos SECOP II integrados.
    \item Requiere definición de hipótesis nulas, modelos (correlación, regresión, panel), y prueba de focalización.
\end{itemize}

\subsection{Análisis espacial condicionado}
\label{sec:espacial_condicionado}

\begin{itemize}
    \item Solo procederá si el diagnóstico de asociación lo justifica.
    \item Requiere: matriz de pesos espaciales (contigüidad, distancia), test de Moran's I, modelos espaciales (SAR, SEM, SDM), detección de clústeres (LISA).
    \item Condicionado a resultados de etapas previas.
\end{itemize}

\subsection{Otras limitaciones técnicas}
\label{sec:otras_limitaciones}

\begin{itemize}
    \item \textbf{Entorno R no versionado:} No hay \texttt{renv.lock} ni \texttt{DESCRIPTION}; versiones de paquetes no fijadas.
    \item \textbf{Fuentes Excel no versionadas:} Archivos en \texttt{datos/} excluidos por \texttt{.gitignore}; reproducibilidad requiere acceso a fuentes originales DANE.
    \item \textbf{Cobertura baja en algunos indicadores:} 28 municipios con cobertura 5PM/pob\_10mas $<0,9$; 26 con menores15 HOGARES $<0,9$; 20 con 60+ HOGARES $<0,95$. Umbrales son parámetros operativos editables en \texttt{R/00\_config.R}, requieren aprobación metodológica.
    \item \textbf{Verificación pendiente 17PM:} El cuadro 17PM está descrito en el índice DANE como ``asistencia escolar'' (igual que 18PM), pero 17PD trata ``nivel educativo''. Si es asistencia, duplica Inasistencia IPM y debe reconsiderarse. Ver \texttt{decisiones\_variables\_IVS\_CNPV2018.csv}, fila 37.
    \item \textbf{Deuda de modularización:} La construcción principal de denominadores e indicadores permanece en gran parte dentro de \texttt{R/02\_base\_final\_indicadores.R} (pendiente extraer a \texttt{R/funciones/}).
\end{itemize}